\chapter{Kết luận và thảo luận}

Khóa luận đã nghiên cứu và phát triển phương pháp luận bộ nhớ sự kiện có bằng chứng nhằm giải quyết bài toán thao tác chuỗi dài cho các mô hình thị giác - ngôn ngữ - hành động (VLA) dưới các ràng buộc khắt khe về tài nguyên tính toán. Thay vì lưu trữ chuỗi quan sát thô gây bùng nổ độ dài ngữ cảnh và phân tán khả năng chú ý của mô hình, phương pháp tiếp cận bài toán qua việc trừu tượng hóa lịch sử tương tác thành các sự kiện chuyển đổi trạng thái thực thể có căn cứ xác thực. Khung phương pháp chứng minh các tính chất toán học bảo đảm tính cập nhật nhân quả không nhìn trước, đồng thời chiếu thông tin tiến độ thao tác vào không gian ngôn ngữ một cách súc tích dưới ngân sách token nghiêm ngặt.

Quá trình đánh giá phân tách trên hai kiến trúc mô hình nền bộc lộ hai phương thức thất bại mang bản chất đối lập khi tiếp nhận ngữ cảnh bộ nhớ. Trên mô hình có dung lượng tham số nhỏ (SmolVLA), việc bổ sung thông tin bộ nhớ vào lời nhắc ngôn ngữ gây xáo trộn phân phối đầu vào, khiến robot suy giảm khả năng phối hợp vận động cơ học ngay từ các mốc tiếp cận vật thể đầu tiên. Trái lại, trên mô hình có năng lực thao tác vật lý vượt trội (X-VLA), robot hoàn thành xuất sắc các mốc chu kỳ kẹp nhấc trung gian nhưng vẫn thất bại ở kết quả toàn phần do tiếp tục lặp lại các động tác thừa mà không tự nhận biết được điểm kết thúc nhiệm vụ. Kết quả thực nghiệm này khẳng định tình trạng mù tiến độ của các mô hình điều khiển liên tục: việc chỉ bổ sung thông tin trạng thái vào lời nhắc ngôn ngữ là chưa đủ để bảo đảm hoàn thành chuỗi dài trọn vẹn.

Bài học kiến trúc then chốt được rút ra là sự cần thiết phải phân tách rạch ròi giữa năng lực sinh hành động của mô hình nền và cơ chế quyết định điểm dừng tác vụ. Để chuyển hóa tỷ lệ hoàn thành mốc cao thành thành công thực tế, hệ thống điều khiển cần được trang bị một cơ chế ngắt chính sách ngoại vi hoạt động song song trong thời gian thực. Khi đó, bằng chứng tích lũy từ bộ nhớ sự kiện kết hợp với cảm nhận trạng thái cánh tay sẽ đóng vai trò là điều kiện dừng độc lập, ngắt luồng điều khiển khi nhiệm vụ đạt trạng thái đích mà không làm xáo trộn phân phối học của mạng nơ-ron điều khiển. Hiệu quả của việc phân tách này đã được chứng thực thực nghiệm khi cơ chế dừng ngoại vi chuyển hóa 9 trong 16 lượt hoàn thành trên E2 và 6 trong 8 lượt trên tập trạng thái seen thành công toàn phần mà không gây bất kỳ lượt dừng sớm nào trên tập seen.

Bên cạnh các phát hiện đạt được, nghiên cứu vẫn ghi nhận một số giới hạn kỹ thuật cần khắc phục. Toàn bộ các thử nghiệm vòng kín hiện mới dừng lại trong môi trường mô phỏng vật lý robosuite và chuẩn LIBERO, chưa phản ánh đầy đủ các biến thiên ánh sáng phức tạp, nhiễu đo lường cảm biến và độ trễ truyền thông trên phần cứng thực tế. Ngoài ra, mô-đun nhận diện thị giác mở dùng để trích xuất sự kiện còn đòi hỏi chi phí tính toán lớn, làm giảm tần số điều khiển khi xử lý đồng bộ trên cùng một thiết bị phần cứng. Đồng thời, các kiểm chứng cơ chế dừng hiện vẫn chịu các giới hạn thực nghiệm: kích thước mẫu còn nhỏ ($40 + 12$ lượt thử nghiệm), mốc không dừng đạt tỷ lệ thành công bằng 0 nên chưa đo được tác hại tiềm tàng trên các lượt chính sách có khả năng tự thành công ($H_0 = 0$), các ngưỡng bản thể vẫn mang tính chất tĩnh, và trần oracle chỉ đạt $10/40$ do kiểm tra dừng cố định tại ranh giới truy vấn chu kỳ 10 bước.

Từ những kết luận và rào cản trên, hướng nghiên cứu tiếp theo sẽ tập trung vào việc tối ưu hóa cơ chế dừng ngoại vi: kiểm tra điều kiện ngắt ở từng bước hành động thay vì chỉ ở ranh giới chu kỳ 10 bước nhằm bắt kịp các khoảng hoàn thành ngắn, và đánh giá trên các chính sách có tỷ lệ tự thành công cơ sở ($H_0 > 0$) để đo lường chính xác tác hại trên các lượt tự hoàn thành. Bên cạnh đó, nghiên cứu cần mở rộng kiểm chứng trên toàn bộ tập chia seen split và các seed ngẫu nhiên khác. Song song với đó, hệ thống cần tối ưu hóa độ trễ của mô-đun trích xuất sự kiện thị giác thông qua các mạng nơ-ron chuyên biệt hoặc chưng cất tri thức sang các mô-đun nhẹ chạy cục bộ để nâng cao tần số điều khiển. Về phạm vi kiểm chứng, nghiên cứu sẽ mở rộng đánh giá vòng kín sang toàn bộ các nhóm tác vụ phụ thuộc lịch sử phức tạp như che khuất và hoán đổi vị trí, hướng tới chuyển giao phương pháp luận bộ nhớ sự kiện lên cánh tay robot thực tế để khẳng định độ tin cậy trong môi trường thực địa.
